We next record two consequences of the simple-graph count that we will
use in the proof for multigraphs. The first finds a marked induced
five-cycle when the clique is large enough. The second bounds the
degrees available around that cycle.

Let $G$ be a simple graph with maximum degree at most a positive
integer $D$, and let $m>0$ of its edges form a strong clique.
Mark these $m$ edges. Write $t=\frac{D}{2}$ and
$\mu=\frac{5}{4} D^2-m$. By \textbf{Lemma~\ref{prop:quadratic-bound}},
$\mu\geq0$. When $D$ is odd, $f(D)=5t^2-t+\frac{1}{4}$. A counterexample
with $m>f(D)$ therefore satisfies
\begin{equation}\label{eq:odd-counterexample}
 \mu\leq t-\frac{5}{4}.
\end{equation}

By \textbf{Lemma~\ref{lem:regularization}}, there is a $D$-regular
simple graph $\widehat G$ containing $G$ as an induced subgraph.
Only the original clique edges are marked. In particular every vertex
incident with a marked edge belongs to $G$.
\textbf{Lemma~\ref{lem:finite-encoding}} supplies the normalized
counts for $\widehat G$.

A marked induced five-cycle is an induced cycle of length five in
$\widehat G$ whose five edges are marked. Let $q_C$ be the normalized
count for the pattern with marked cycle $0,1,2,3,4,0$ and root $01$.
Every color-preserving map of this pattern is injective. Adjacent
positions have different images, and given two nonadjacent positions,
there is a third position adjacent to exactly one of them.
In the balanced blow-up with $b$ vertices in each part and all
edges marked, there are $10b^5$ such maps, from the ten cyclic
orderings and the choice of one vertex per part. Hence
$q_C=\frac{10b^5}{2(5b^2)(2b)^3}=\frac{1}{8}$ in that example.

Fix an ordered image $c_0,\ldots,c_4$ of the cycle, taking indices
modulo five. Let $U_i$ consist of the vertices whose neighbors on
the cycle are exactly $c_{i-1},c_{i+1}$ and whose edges to both
cycle vertices are marked. The sets $U_i$ are pairwise disjoint,
contain $c_i$ in $U_i$, and lie in $V(G)$ since every vertex in
them is incident with a marked edge.
\textbf{Figure~\ref{fig:odd-cycle-parts}} illustrates $U_0$. Define
\begin{equation}\label{eq:odd-core-defs}
 a_i=|U_i|,\qquad B=\bigcup_{i=0}^4 U_i,\qquad
 s_i=D-a_{i-1}-a_{i+1},\qquad R=\sum_{i=0}^4 s_i.
\end{equation}
By the definitions of the five sets, the neighbors of $c_i$ in $B$
are exactly $U_{i-1}\cup U_{i+1}$. Since $\widehat G$ is $D$-regular,
$s_i=|N_{\widehat G}(c_i)\setminus B|$. Summing over the five cycle
vertices gives $R=\sum_i|N_{\widehat G}(c_i)\setminus B|$, an upper
bound on the number of edges of $G$ from the cycle to $V(G)\setminus B$.

\begin{figure}[htbp]
\centering
\begin{tikzpicture}[every node/.style={font=\small}]
  \node[graph vertex,label=below:$c_0$] (c0) at (0,1.4) {};
  \node[graph vertex,label=left:$c_1$] (c1) at (-1.55,0.45) {};
  \node[graph vertex,label=below:$c_2$] (c2) at (-0.95,-1) {};
  \node[graph vertex,label=below:$c_3$] (c3) at (0.95,-1) {};
  \node[graph vertex,label=right:$c_4$] (c4) at (1.55,0.45) {};
  \node[graph vertex,label=above:$x$] (x) at (0,2.35) {};
  \draw[marked edge] (c0)--(c1)--(c2)--(c3)--(c4)--(c0);
  \draw[marked edge] (c1)--(x)--(c4);
\end{tikzpicture}
\caption{A vertex $x$ in $U_0$ and its two marked edges to the cycle.}
\label{fig:odd-cycle-parts}
\end{figure}


To express $R$ in terms of the counts, let $L_R$ be the following
linear form in the six-position coordinates. The restriction to the
first five positions is the marked induced cycle. The coefficient
of a labeled extension is the number of neighbors of the sixth
position on the cycle if that position either has an unmarked edge
to the cycle or is adjacent to two consecutive cycle positions.
All other extensions have coefficient zero. A vertex outside $B$
with a cycle neighbor satisfies at least one of these two tests.
Indeed, if it fails both, it has at most two cycle neighbors and
all edges to them are marked. With only one such neighbor $c_i$,
its marked edge to $c_i$ has no endpoint in common with the
opposite cycle edge $c_{i+2}c_{i+3}$ and no edge joins their endpoints.
With two nonconsecutive marked neighbors, it belongs to one of
the sets $U_i$. Conversely, vertices of $B$ fail both tests.

Thus the sum over the sixth position is $R$ for each ordered
cycle. The denominators in \eqref{eq:q-definition} for five and
six positions differ by a factor of $D$. If $q_C>0$, this gives
$L_R=\frac{q_C}{D}\mathbb E[R]$, where the expectation is over a
uniformly chosen color-preserving map of the cycle pattern. If $q_C=0$,
there is no such map and $L_R=0$.

The five-cycle certificates give
\begin{equation}\label{eq:odd-moment-bounds}
 \frac{\mu}{D^2}\geq\gamma_C\Bigl(\frac{1}{8}-q_C\Bigr)-\epsilon_C,
 \qquad
 \frac{\mu}{D^2}\geq\gamma_R L_R-\epsilon_R.
\end{equation}
Both errors $\epsilon_C$ and $\epsilon_R$ are nonnegative. For the
first inequality we use either
$\gamma_C\geq\frac{99999}{100000}$ and $\epsilon_C\leq2\cdot10^{-9}$,
or $\gamma_C\geq\frac{6449}{10000}$ and $\epsilon_C\leq10^{-9}$.
For the second we use $\gamma_R\geq\frac{56425}{100000}$ and
$\epsilon_R\leq10^{-9}$. The files and checks are specified in
\cite[Sections~B and~C]{FGSTSupplement}.

Suppose that $\mu\leq\mu_*$, $\gamma_C\geq g_C>0$ and
$\gamma_R\geq g_R>0$. Put
$q_*=\frac{1}{8}-\frac{\mu_*+\epsilon_C D^2}{g_C D^2}$.
The first inequality in \eqref{eq:odd-moment-bounds} gives
$q_C\geq q_*$ when $q_C<\frac{1}{8}$, and the bound holds directly
when $q_C\geq\frac{1}{8}$ because $q_*\leq\frac{1}{8}$. If $q_*>0$,
using the second inequality and $L_R=\frac{q_C}{D}\mathbb E[R]$
gives a cycle with $R\leq\mathbb E[R]$ and hence
\begin{equation}\label{eq:odd-choose-anchor}
 R\leq\frac{\mu_*+\epsilon_R D^2}{g_R Dq_*}.
\end{equation}
Verified upper bounds on either error can be used in this estimate.

\begin{lemma}\label{lem:simple-cycle-counts}
For the sets defined in \eqref{eq:odd-core-defs}, every vertex of
$U_i$ is adjacent to every vertex of $U_{i+1}$. Moreover,
$|B|=\frac{5D-R}{2}$ and $a_i=\frac{D-R}{2}+s_{i-2}+s_{i+2}$ for every $i$.
\end{lemma}

\begin{proof}
For every $x\in U_i$ and $y\in U_{i+1}$ we have $xy\in E(G)$. The
marked edges $xc_{i-1}$ and $yc_{i+2}$ are disjoint. The
definitions of the parts exclude $xc_{i+2}$ and $yc_{i-1}$, and
$c_{i-1}c_{i+2}$ is excluded because the cycle is induced. Since
the marked edges form a clique in $L(G)^2$, we must have
$xy\in E(G)$. Thus a vertex of $U_i$ has at most $s_i$ neighbors
outside $U_{i-1}\cup U_{i+1}$.

The sum $A=\sum_i a_i s_i$ bounds the number of incidences at vertices
of $B$ after excluding the edges between $U_i$ and $U_{i+1}$ for all
$i$. Each remaining edge with both ends in $B$ contributes two
incidences. Put $e_0=\sum_i a_i a_{i+1}$, the number of edges between
parts $U_i$ and $U_{i+1}$.
Solving \eqref{eq:odd-core-defs} gives
\begin{equation}\label{eq:odd-core-sums}
 a_i=t-\frac{R}{2}+s_{i-2}+s_{i+2},\qquad
 |B|=\frac{5D-R}{2}.
\end{equation}
Summing the five equations for $s_i$ gives the formula for $|B|$.
For $a_i$, the two disjoint pairs of indices
$\{i-2,i+1\}$ and $\{i-1,i+2\}$ are the neighbor pairs of
$i+2$ and $i-2$, respectively. Subtracting their sizes from $|B|$
gives $a_i=|B|-2D+s_{i-2}+s_{i+2}$.
\end{proof}

By substituting \eqref{eq:odd-core-sums} into $\sum_i a_i s_i$
and using $2e_0=\sum_i a_i(D-s_i)$, we obtain
\begin{equation}\label{eq:odd-core-capacity}
 A=tR-\frac{R^2}{2}+2\sum_i s_i s_{i+2},
 \qquad
 e_0=5t^2-\frac{Rt+A}{2}.
\end{equation}
When $D$ is odd, $R=5D-2|B|$ is odd as well. To bound $\sum_i s_i s_{i+2}$,
consider the cycle on the indices with edges $i,i+2$. If two
nonadjacent vertices have positive weights, move their combined
weight onto the vertex with the higher sum of neighboring weights.
With the other weights fixed, the sum of products does not decrease.
Repeating leaves positive weights on a clique of size at most two.
The sum of products is then at most $\frac{R^2}{4}$. When $D$ is odd, integrality reduces this bound to $\frac{R^2-1}{4}$, so
\eqref{eq:odd-core-capacity} gives $A\leq tR-\frac{1}{2}$.

\begin{lemma}\label{lem:odd-unit-slack}
If an induced five-cycle with five marked edges has $R=1$, then $D$ is odd and $m\leq f(D)$.
\end{lemma}

\begin{proof}
Let $i$ be the unique index with $s_i=1$. Every vertex of the other
four parts has degree $D$ from the consecutive pairs alone. Thus a
marked edge outside the consecutive pairs has both ends in
$U_i\cup(V(G)\setminus B)$. The marked edge $c_{i+2}c_{i+3}$ has both
ends in parts whose consecutive neighbors are outside $U_i$. Both
vertices already have degree $D$ from those pairs, so neither is
an endpoint of, or adjacent to an endpoint of, the proposed marked
edge. This contradicts the clique property. Hence $m\leq e_0$.
For $R=1$, \eqref{eq:odd-core-capacity} gives
$A=t-\frac{1}{2}$ and $e_0=5t^2-t+\frac{1}{4}=f(D)$.
\end{proof}
